% Kandydat restrukturyzacji; nie jest aktywnym rozdziałem Atlasu.
\section{Zbieżność jednostajna szeregów Fouriera}
\providecommand{\atlasref}[2]{#2}

\begin{lemat}[Zanikanie współczynników gładkiej funkcji]
\label{lem:dt-series-decay}
Jeśli $f\in C^m(\mathbb T)$, $m\geq1$ całkowite, to dla $k\ne0$
$|\widehat f(k)|\le\|f^{(m)}\|_{L^1(-\pi,\pi)}/(2\pi|k|^m)$.
Jeśli $f\in C^\infty(\mathbb T)$, to szereg Fouriera wraz
ze wszystkimi szeregami otrzymanymi przez różniczkowanie
zbiega jednostajnie.
\end{lemat}
\begin{proof}
Dla $k\ne0$ całkowanie przez części daje
\[
 \widehat f(k)=\frac1{ik}\widehat{f'}(k),
\]
ponieważ $f(\pi)e^{-ik\pi}=f(-\pi)e^{ik\pi}$.
Powtarzając ten rachunek $m$ razy (pochodne
do rzędu $m-1$ mają zgodne wartości na końcach
przedziału), dostajemy
$\widehat f(k)=(ik)^{-m}\widehat{f^{(m)}}(k)$,
co po oszacowaniu modułu całki daje żądaną nierówność.

Dla $f\in C^\infty$ i ustalonego całkowitego
$r\geq0$ wybieramy $m\geq r+2$. Wówczas
\[
 \sum_{k\ne0}|k|^r|\widehat f(k)|
 \leq\frac{\|f^{(m)}\|_{L^1(-\pi,\pi)}}{2\pi}
 \sum_{k\ne0}|k|^{r-m}<\infty.
\]
Kryterium Weierstrassa daje bezwzględną
i jednostajną zbieżność szeregu
$\sum_{k\in\mathbb Z}(ik)^r\widehat f(k)e^{ikx}$,
przy czym dla $r=0$ wyraz $k=0$ oznacza
$\widehat f(0)$, a dla $r>0$ wynosi zero.
Oznaczmy jego ciągłą sumę przez $G_r$.
Z całkowania przez części także dla $k=0$
(jeśli $r>0$, całka z $f^{(r)}$ po okresie
jest zerowa) współczynniki tego szeregu
są dokładnie współczynnikami Fouriera $f^{(r)}$.
Twierdzenie~\ref{thm:dt-series-parseval} daje więc
zbieżność tych samych sum częściowych do $f^{(r)}$
w normie $L^2$. Z jednostajnej zbieżności
mamy zarazem zbieżność w $L^2$ do $G_r$.
Jednoznaczność granicy w $L^2$ daje
$G_r=f^{(r)}$ prawie wszędzie, a ponieważ
obie funkcje są ciągłe, równość zachodzi
w każdym punkcie: niezerowa różnica
w jednym punkcie byłaby niezerowa na
pewnym otwartym łuku dodatniej miary.
W szczególności $G_0=f$ i wszystkie
różniczkowania wyraz po wyrazie prowadzą
do właściwych pochodnych.
\end{proof}
